\section{To Be of the Right}

\begin{quotex}
This essay by \textbf{Julius Evola} was originally published in the journal ``Roma'' on 19 March 1973 under the title ``Essere di Destra''.
\end{quotex}

Right and Left are designations that refer to a political system already in crisis. In traditional systems of government, they were non-existent, at least if taken in their current meaning. In them there could be an opposition party, but not a revolutionary one, nor one that put the system itself in question, but rather it was loyalist and, in a certain way, workable. So in England, where one could speak of a ``very loyal opposition of His Majesty''. Things changed after facing the subversive movements of more recent times, and it is known that at their origin, the Right and the Left were defined on the basis of the places the opposed parties respectively occupied in parliament.

Depending on the level, the Right assumes distinct meanings. There is an economic Right based on capitalism, not without legitimacy if it does not abuse its power and if its antithesis is socialism and Marxism.

As to a political Right, it strictly acquires its full significance if a monarchy exists in an organic State: as was the case especially in central Europe, and also partly in conservative England.

But we can also disregard institutional presuppositions and speak of a Right in terms of a spiritual orientation and vision of the world. Then to be of the Right means, beyond being against democracy and every ``social'' mythology, to defend the values of Tradition as spiritual, aristocratic, and warrior values (in a derivative way, also with reference to a strict military tradition as, for example, happened in Prussianism). It means more than nourishing a certain disdain for intellectualism and for the bourgeois fetishism of the ``cultured man''. (The representative of an old Piedmontese family paradoxically had this to say: ``I divide our world into two classes: the nobility and those who have a university degree'' and Ernst Jünger, in support, valorized the antidote constituted by a ``healthy analphabetism'').

To be of the Right means also to be conservative, although not in a static sense. The obvious presumption is that there is something of substance worthy of being conserved, which however puts us facing a difficult problem. In reference to what constituted the immediate past of Italy after its unification: the Italy of the nineteenth century has certainly not left us a heritage of superior values to preserve or deeds to serve as the base. Even going further back in Italian history, we encounter only sporadic positions of the Right; it lacked a formative unitary force which existed in other nations, solidified over time from ancient monarchical traditions of an aristocratic oligarchy.

However, in affirming that a Right must not be characterized by a static conservatism, we mean that there must be certain values or certain ideas based on solid ground, but that different expressions must be given to them, adequate to the development of the times, in order to not let them be bypassed, in order to take back, control, and incorporate everything that little by little is manifested in changing situations. This is the only sense in which a man of the Right can conceive ``progress''; it is not a simple movement forward, as too many often think, especially among the left; Bernanos was able to speak aptly of a ``escape forward'' in this contest (``imbeciles, where are you fleeing ahead?''). ``Progressivism'' is a foible foreign to every positions of the Right. It is also because in a general consideration of the course of history --- with reference to spiritual values rather than to material values, technical achievements, etc. --- the man of the Right only recognized a descent, not progress nor a true ascent. The developments of current society can only confirm this conviction.

The positions of a Right are necessarily anti-corporate, anti-plebeian, and aristocratic; thus their positive counterpart will see value in the affirmation of the ideal of a well-structured, organic, hierarchical State, directed from a principle of authority. In this last regard, they therefore overlook the difficulties with regards to that from which such a principle can draw its foundation and its blessing. It is obvious that it cannot come from below, from the \emph{demos}, in which, without offending the Mazzinians of yesterday and today, it does not express in the least the \emph{vox Dei} voice of God , and if anything, the contrary. And one must also exclude the dictatorial and Bonapartist solutions, which can only have transient value, in emergency situations with contingent and short-term goals.

Again, we find ourselves compelled to refer instead to dynastic continuity, provided we keep in view at least what was called ``authoritarian constitutionalism'' in monarchic regimes, that is, a power that is not purely representative, but also active and regulatory. This is on the level of that ``decisionism'' which De Maistre and Donoso Cortes had previously spoken about, in reference to decisions constituting the last resort, with all the responsibility tied to it and which must be assumed in person, when one is found facing the necessity of a direct intervention because the existing order has entered a crisis or new forces push onto the political scene.

Let us repeat however that the rejection in these terms of a ``static conservatism'' does not concern the level of principles. For the man of the Right, principles always constitute the solid base, the terra firma in the face of change and contingencies, and there the counter-revolution deserves a precise catchword. If we want, we can refer here instead to the formula, paradoxical only in appearance, of a ``conservative revolution''. It concerns all the initiatives that are imposed through the removal of negative factual situations, necessary for a restoration and for an adequate revival of what has an intrinsic value and cannot be an object of discussion. In effect, in conditions of crisis and subversion, it can be said that nothing has a character so revolutionary as much as the renewal of such values. An ancient saying \emph{usu vetera novant} the use of old and new , and it emphasizes the same context: the renovation that can actualize the revival of the ``ancient'', i.e., the unchangeable traditional legacy.

With this, we believe that the positions of the man of the Rights are sufficiently clarified.


\flrightit{Posted on 2012-10-29 by Cologero}

\begin{center}* * *\end{center}

\begin{footnotesize}
\begin{sffamily}
\texttt{Logres on 2012-10-30 at 20:33 said:}

I corresponded briefly with John Michael Greer (Archdruid of North America) over a comment he made re: Evola, and here is what he said (reposted with permission)-

Question -- ``Would be interested in your take on Evola's ``garbling'' of celestial/telluric currents... ''

Greer:

``I doubt it will be congenial to the Gornahoor audience. Druidry shares much in common with traditional Taoism --- I can say this with some basis in experience, having studied with a Taoist master and put a decade into one of the old temple styles of t'ai chi ch'uan --- and the sense of the nature of polarity is an important piece of this common ground. The notion that any polarity ought to be resolved by the conquest of one extreme by the other , in both traditions, is simply wrong --- and this is all the more true as Evola's hypothetical Uranian and Demetrian Traditions are at best, dubiously historic reifications of eternal spiritual realities. I'd point out that in alchemy, the King does not rape and enslave the Queen --- not unless you want your alchemical working to misfire from the very beginning and produce a poison rather than an elixir. Instead, the King and Queen marry, mate, and fuse into the twofold Rebis, from which the living Stone is born. In the same way, yin and yang, the solar and telluric currents, the Uranian and Demetrian principles are incomplete and unbalanced when separated from one another, and produce not only balance and harmonious change but a new, third thing when brought together in the right way. Evola's problem, as I see it, was that he was too heavily influenced by the popular intellectual culture of his own time --- his ranking of warriors above priests, for example, was taken straight out of Nietzsche, and the whole Uranian/Demetrian dualism is a projection of ideas borrowed from Weininger's Geschlecht und Charakter --- most of the radical Right in his time was influenced by Weininger, which makes it all the more curious that nobody reads Weininger at all these days. It's all too common for those who believe they've found a timeless tradition to define it in terms that will look profoundly dated a generation or two thereafter!''

I'll note that Cologero has read Weiniger, and that it is in the library. Although I find his comments useful up to a point, it seems to me to miss the point: Evola was a subtle and complex character addressing a specific time and danger -- Greer seems to be engaged in polemics here, just like a typical modern. All the same, if anyone has pointers as to what to write back to him, I would be much interested... 

\hfill

\texttt{Cologero on 2012-10-31 at 01:14 said:}

I don't know what you want to say to him, Logres, since we have often criticized Evola here on similar grounds. Although Evola himself was versed in Taoism, even publishing an Italian translation of the Tao Te Ching, his duality of Solar and Lunar as Good and Evil is off the mark. He was too influenced by the likes of Bachofen and Weininger. It is one thing to use them to illustrate points, but quite another to inject them into his system. Part of Evola's ``complexity'' is the result of his prejudices. Evola should always be read in the light of Guenon, whom Evola acknowledged as the Master of the 20th century.

As we have amply documented, to his dying day, Evola preferred his own philosophical system to proper metaphysics. I am aware that that makes him more ``exciting'' to certain casts of mind, but to read Evola without a firm foundation in Guenon (or even one's own traditional teachings) will lead to serious misunderstandings.

\hfill

\texttt{Avery M on 2012-10-31 at 09:30 said:}

Cologero, if you were to expand that single comment into a pamphlet or an ebook, it would make a fine introduction to reading Guenon and Evola side by side. One can eventually glean this stuff on one's own, especially if one tries to struggle through Evola's monstrous \emph{Metaphysics of Sex} or Weiniger, but I think people would prefer a guide so that they don't have to do that.

\hfill

\texttt{Logres on 2012-10-31 at 15:33 said:}

Agree with Avery. As to reply, agreed also that Evola ignores the feminine for purposes of constructing his world (that's clear enough from Metaphysics of Sex, one of the few Evola books I've read cover to cover). The problem is that he's ``constructing'' at all -- however, his emphasis on Monotheism \& the Empire is something the Druidry revival would eschew; and he doesn't so much ``garble'' celestial and teluric currents as separate them arbitrarily in order to purge what he doesn't have a use for. But this is exactly what Druidry does with paganism and monotheism. Or so it seems to me... 

\hfill

\texttt{Saladin on 2012-10-31 at 19:57 said:}

I am not sure if one MUST ``read Evola in the light of Guenon''. Just as there are many different manifestations of the Primordial Tradition, as religions, there can be different variations in interpreting that same Sanatana Dharma! Personally, I am indebted to both of them but I can see why Guenon should be given precedence.

As for replying to Mr. Greer: Please make clear to him that there exists two fundamentally different Taoisms. The Traditional one is the Taoism of Lao Tze and Chuang Tze (which in our era is almost impossible to find in any organized or initiatory form) but the other Taoism is nothing but a degenerate form of Shamanism. What is known by Taoism in Taiwan (and in mainland China if any remnants have survived the vicissitudes of Communism) by the term Taoism is a disparate practices of third rate Magic and popular superstition. ``It is a set of ancient religions, ancestor worship, animism, shamanism, and magic, mixed to varying degrees with Buddhism, Confucianism, and Traditional Taoism.'' It appears what Mr. Greer refers to is the latter.

\hfill

\texttt{aperion on 2012-10-31 at 20:03 said:}

No Avery, that would be too easy, nay, a a cop out. I mean is it really that difficult for modern youth to read a Weiniger (tempted to tap a double ``g'' but won't) through and through or even a very comprehensible ``Metaphysics of Sex''? To hell with guides, both books are entirely coherent to anyone outside of specialist academic categories for those who care to actually read and comprehend books. Basically you want a summery or even care i day a Sacrament(!)- why do we even bother some my cry- , so you can avoid the pain of study. My diagnosis: man up beotch.

\hfill

\texttt{Logres on 2012-10-31 at 21:17 said:}

In fairness to him, he states he reads the original texts (maybe in translation?) -- the man has a first rate mind and is a gentlemen.


\hfill

\end{sffamily}
\end{footnotesize}

